\section{The Gettier Problem as a Structural Diagnostic}

\subsection{Overview}

The Gettier problem is usually treated as a problem in the definition of knowledge. This paper reads it differently. It is a structural diagnostic showing that a stopping structure cannot guarantee a dynamical property.

The classical definition of knowledge as justified true belief attempts to stabilize knowledge by adding conditions. Gettier cases show that the stabilization fails. The failure is not accidental. It occurs because knowledge belongs to the layer of stopping structures, whereas intelligence belongs to ongoing dynamics.

\subsection{Knowledge as stopping structure}

Knowledge is useful because it stops. A proposition can be stored, cited, transmitted, and reused only because it has become sufficiently stable within an observational window.

This stability is also the source of the problem. A knowledge state can be true, believed, and justified while still being dynamically fragile. It may fail under counterfactual variation, under later revision, or under a wider inferential context.

The point is not that knowledge is useless. The point is that knowledge cannot carry the full dynamical property of intelligence.

\subsection{Why Gettier cases recur}

A Gettier case arises when a stopped state satisfies available criteria while failing to preserve the motion that those criteria were supposed to certify. Each repair adds a further stopping condition. Each further condition creates a new surface at which the dynamic can fail again.

Thus the Gettier problem is not merely a collection of counterexamples. It reveals a mismatch:

\begin{quote}
A stopping structure is being asked to guarantee an ongoing motion.
\end{quote}

No refinement of the stopping structure can remove the mismatch completely, because the property being sought is not located at the knowledge layer.

\subsection{Relation to the four-layer decomposition}

In the decomposition of \S{}2, knowledge is the layer of stability, whereas intelligence is the layer of revisable motion. Gettier cases occur when the stability of knowledge is mistaken for the motion of intelligence.

This also explains why knowledge can be necessary without being sufficient. Intelligence generates and uses knowledge, but it is not identical with knowledge. A system may possess many stable stopping structures and still fail to sustain intelligent dynamics.

\subsection{Implications for later sections}

The Gettier problem prepares the argument of \S{}7 and \S{}8. If correctness cannot stabilize intelligence, then accumulated correctness may not only fail to help; it may intensify instability. \S{}7 shows this dynamically through branching explosion. \S{}8 shows the structural reason: the stopping operator is not internal to the individual inferential system.

\subsection{Summary}

\begin{quote}
The Gettier problem is not primarily a defect in the definition of knowledge. It is a signal that knowledge and intelligence occupy different structural layers. Knowledge stops. Intelligence revises.
\end{quote}
